% GENERATED FILE. Edit canonical lessons, then run npm run generate:books.
% canonical-source-hash: fnv1a64:9709ee271a923382
% canonical-lessons: HI-C368-sudharnaa, HI-C368-vikas, HI-C368-unnati, HI-C368-anusasan, HI-C368-nikhaarna

\chapter{Making Things Better}
\label{ch:hi-making-things-better}
\hlchaptermodality{368}

\begin{chapteropening}
\textbf{By the end of this chapter:} \emph{I can say how a place or a person improved.}

Complete the last lesson of chapter 368: I can say how a place or a person improved.
\end{chapteropening}

\section[\texorpdfstring{\hi{सुधरना} (sudharnā)}{सुधरना (sudharnā)}]{\hi{सुधरना} (sudharnā) — to get better, to mend one's ways}

\label{lesson:HI-C368-sudharnaa}



\begin{quote}
Before the new one: say the Hindi for criticism, then the Hindi for to endure.
\end{quote}

\subsection*{You'll want to know: \hi{सुधरना}}
\textbf{\hi{सुधरना}} — \emph{sudharnā} — \textquotedblleft{}to get better, to mend one's ways\textquotedblright{}.

\begin{quote}
\hi{दवा} \hi{के} \hi{बाद} \hi{उसकी} \hi{तबीयत} \hi{सुधर} \hi{रही} \hi{है।}
\end{quote}

\begin{quote}
\emph{davā ke bād uskī tabīyat sudhar rahī hai.}
\end{quote}

\begin{quote}
After the medicine his health is getting better.
\end{quote}

Past (he …): \hi{वह} \hi{सुधर} \hi{गया}, \emph{vah sudhar gayā}.

\begin{cousinweb}
\textbf{\hi{सुधरना}} (\emph{sudharnā}) is inherited from \emph{*śuddhakāra}, \textquotedblleft{}making clean\textquotedblright{}, a form CDIAL reconstructs from its descendants (CDIAL 12521).
\end{cousinweb}

\begin{grammarlens}[title={What you've built}]
One more word for giving reasons and opinions.
\end{grammarlens}

\subsection*{Your turn}
\begin{itemize}
\raggedright
  \item \emph{Say it:} \emph{sudharnā}
  \item \emph{Say it:} \emph{sudharnā}, once more
  \item \emph{Say it:} \emph{davā ke bād uskī tabīyat sudhar rahī hai.} — After the medicine his health is getting better.
  \item \emph{Recall:} say the Hindi for criticism, then the Hindi for to endure, then say \emph{sudharnā} again
\end{itemize}

\begin{tcolorbox}[breakable,colback=teal!4,colframe=teal!35!black,title={Before you move on}]
What does \hi{सुधरना} mean? (\textquotedblleft{}To get better, to mend one's ways\textquotedblright{}.) Say it once more.
\end{tcolorbox}
\section[\texorpdfstring{\hi{विकास} (vikās)}{विकास (vikās)}]{\hi{विकास} (vikās) — development}

\label{lesson:HI-C368-vikas}



\begin{quote}
Before the new one: say the Hindi for to endure, then the Hindi for to get better.
\end{quote}

\subsection*{You'll want to know: \hi{विकास}}
\textbf{\hi{विकास}} — \emph{vikās} — \textquotedblleft{}development\textquotedblright{}.

\begin{quote}
\hi{गाँव} \hi{के} \hi{विकास} \hi{के} \hi{लिए} \hi{सड़क} \hi{ज़रूरी} \hi{है।}
\end{quote}

\begin{quote}
\emph{gā̃v ke vikās ke liye saṛak zarūrī hai.}
\end{quote}

\begin{quote}
A road is necessary for the village's development.
\end{quote}

Masculine.

\begin{cousinweb}
\textbf{\hi{विकास}} (\emph{vikās}) is a tatsama, borrowed from Sanskrit \emph{vikāsa}, \textquotedblleft{}opening, expanding\textquotedblright{} (Monier-Williams).
\end{cousinweb}

\begin{grammarlens}[title={What you've built}]
One more word for giving reasons and opinions.
\end{grammarlens}

\subsection*{Your turn}
\begin{itemize}
\raggedright
  \item \emph{Say it:} \emph{vikās}
  \item \emph{Say it:} \emph{vikās}, once more
  \item \emph{Say it:} \emph{gā̃v ke vikās ke liye saṛak zarūrī hai.} — A road is necessary for the village's development.
  \item \emph{Recall:} say the Hindi for to endure, then the Hindi for to get better, then say \emph{vikās} again
\end{itemize}

\begin{tcolorbox}[breakable,colback=teal!4,colframe=teal!35!black,title={Before you move on}]
What does \hi{विकास} mean? (\textquotedblleft{}Development\textquotedblright{}.) Say it once more.
\end{tcolorbox}
\section[\texorpdfstring{\hi{उन्नति} (unnati)}{उन्नति (unnati)}]{\hi{उन्नति} (unnati) — progress, advancement}

\label{lesson:HI-C368-unnati}



\begin{quote}
Before the new one: say the Hindi for to get better, then the Hindi for development.
\end{quote}

\subsection*{You'll want to know: \hi{उन्नति}}
\textbf{\hi{उन्नति}} — \emph{unnati} — \textquotedblleft{}progress, advancement\textquotedblright{}.

\begin{quote}
\hi{पढ़ाई} \hi{से} \hi{देश} \hi{की} \hi{उन्नति} \hi{होती} \hi{है।}
\end{quote}

\begin{quote}
\emph{paṛhāī se deś kī unnati hotī hai.}
\end{quote}

\begin{quote}
A country advances through education.
\end{quote}

Feminine.

\begin{cousinweb}
\textbf{\hi{उन्नति}} (\emph{unnati}) is a tatsama, borrowed from Sanskrit \emph{unnati}, \textquotedblleft{}rising; advancement\textquotedblright{} (Monier-Williams).
\end{cousinweb}

\begin{grammarlens}[title={What you've built}]
One more word for giving reasons and opinions.
\end{grammarlens}

\subsection*{Your turn}
\begin{itemize}
\raggedright
  \item \emph{Say it:} \emph{unnati}
  \item \emph{Say it:} \emph{unnati}, once more
  \item \emph{Say it:} \emph{paṛhāī se deś kī unnati hotī hai.} — A country advances through education.
  \item \emph{Recall:} say the Hindi for to get better, then the Hindi for development, then say \emph{unnati} again
\end{itemize}

\begin{tcolorbox}[breakable,colback=teal!4,colframe=teal!35!black,title={Before you move on}]
What does \hi{उन्नति} mean? (\textquotedblleft{}Progress, advancement\textquotedblright{}.) Say it once more.
\end{tcolorbox}
\section[\texorpdfstring{\hi{अनुशासन} (anuśāsan)}{अनुशासन (anuśāsan)}]{\hi{अनुशासन} (anuśāsan) — discipline}

\label{lesson:HI-C368-anusasan}



\begin{quote}
Before the new one: say the Hindi for development, then the Hindi for advancement.
\end{quote}

\subsection*{You'll want to know: \hi{अनुशासन}}
\textbf{\hi{अनुशासन}} — \emph{anuśāsan} — \textquotedblleft{}discipline\textquotedblright{}.

\begin{quote}
\hi{खिलाड़ियों} \hi{में} \hi{अनुशासन} \hi{बहुत} \hi{ज़रूरी} \hi{है।}
\end{quote}

\begin{quote}
\emph{khilāṛiyoṁ meṁ anuśāsan bahut zarūrī hai.}
\end{quote}

\begin{quote}
Discipline is very important in players.
\end{quote}

Masculine.

\begin{cousinweb}
\textbf{\hi{अनुशासन}} (\emph{anuśāsan}) is a tatsama, borrowed from Sanskrit \emph{anuśāsana}, \textquotedblleft{}instruction, a precept\textquotedblright{} (Monier-Williams).
\end{cousinweb}

\begin{grammarlens}[title={What you've built}]
One more word for giving reasons and opinions.
\end{grammarlens}

\subsection*{Your turn}
\begin{itemize}
\raggedright
  \item \emph{Say it:} \emph{anuśāsan}
  \item \emph{Say it:} \emph{anuśāsan}, once more
  \item \emph{Say it:} \emph{khilāṛiyoṁ meṁ anuśāsan bahut zarūrī hai.} — Discipline is very important in players.
  \item \emph{Recall:} say the Hindi for development, then the Hindi for advancement, then say \emph{anuśāsan} again
\end{itemize}

\begin{tcolorbox}[breakable,colback=teal!4,colframe=teal!35!black,title={Before you move on}]
What does \hi{अनुशासन} mean? (\textquotedblleft{}Discipline\textquotedblright{}.) Say it once more.
\end{tcolorbox}
\section[\texorpdfstring{\hi{निखारना} (nikhārnā)}{निखारना (nikhārnā)}]{\hi{निखारना} (nikhārnā) — to polish (a skill), to bring out the best in}

\label{lesson:HI-C368-nikhaarna}



\begin{quote}
Before the new one: say the Hindi for advancement, then the Hindi for discipline.
\end{quote}

\subsection*{You'll want to know: \hi{निखारना}}
\textbf{\hi{निखारना}} — \emph{nikhārnā} — \textquotedblleft{}to polish (a skill), to bring out the best in\textquotedblright{}.

\begin{quote}
\hi{अभ्यास} \hi{से} \hi{वह} \hi{अपनी} \hi{हिंदी} \hi{निखार} \hi{रहा} \hi{है।}
\end{quote}

\begin{quote}
\emph{abhyās se vah apnī hindī nikhār rahā hai.}
\end{quote}

\begin{quote}
With practice he is polishing his Hindi.
\end{quote}

Past, with \hi{ने} (he or she …): \hi{उसने} \hi{अपना} \hi{गाना} \hi{निखारा}, \emph{usne apnā gānā nikhārā}.

\begin{cousinweb}
\textbf{\hi{निखारना}} (\emph{nikhārnā}) is inherited from \emph{*niḥkṣārayati}, \textquotedblleft{}causes to flow out\textquotedblright{}, a form CDIAL reconstructs from its descendants (CDIAL 7098). CDIAL glosses the Hindi verb \textquotedblleft{}to clarify, to strain\textquotedblright{}; polishing a skill is not among its glosses, and may be the same idea of clearing.
\end{cousinweb}

\begin{grammarlens}[title={What you've built}]
One more word for giving reasons and opinions.
\end{grammarlens}

\subsection*{Your turn}
\begin{itemize}
\raggedright
  \item \emph{Say it:} \emph{nikhārnā}
  \item \emph{Say it:} \emph{nikhārnā}, once more
  \item \emph{Say it:} \emph{abhyās se vah apnī hindī nikhār rahā hai.} — With practice he is polishing his Hindi.
  \item \emph{Recall:} say the Hindi for advancement, then the Hindi for discipline, then say \emph{nikhārnā} again
\end{itemize}

\begin{tcolorbox}[breakable,colback=teal!4,colframe=teal!35!black,title={Before you move on}]
What does \hi{निखारना} mean? (\textquotedblleft{}To polish (a skill), to bring out the best in\textquotedblright{}.) Say it once more.
\end{tcolorbox}
